\section{可行性与时间评估}

在前文已经完成选题重构、方向比较以及 \texttt{Belief-State Controller} 的问题定义之后，最后必须回答一个最现实的问题：在当前时间已经接近三月底、距离毕业设计收尾不足两个月的条件下，这一方向是否仍然具有可行性。这个问题不仅关系到技术安排，更关系到选题是否真正能够收敛为一个可答辩的毕业设计。因此，本节不再讨论该方向“理论上是否有意义”，而是集中分析它在现阶段是否具备现实落地条件，以及在有限时间内应如何控制范围。

\subsection{可行性的前提：现有系统已经提供了足够的实验平台}

判断当前方向是否可行，首先要看是否仍需从零搭建系统。如果答案是“需要重写一套全新的运行时平台”，那么在现阶段继续转向显然不现实；但如果现有系统已经提供了大部分实验所需的基础，那么控制器方向的可行性就会显著提高。结合前文对当前系统的分析，可以看到后者更接近真实情况。

当前系统已经具备显式有限状态机、阶段化蛋白设计流程、候选生成与质量门控、由 \texttt{retry}、\texttt{patch} 和 \texttt{replan} 组成的恢复链，以及人工介入机制。这意味着工作流的基本执行逻辑、可恢复执行链路和人工确认链路都已经存在。对于 \texttt{Belief-State Controller} 而言，这一点非常关键，因为控制器本身并不需要重新实现整个任务系统，它所依附的正是这些已有的运行时结构。换言之，当前课题不是从零开始做一套 agent 平台，而是在既有平台之上补足一个“如何依据运行时状态做治理决策”的核心算法层。

更重要的是，系统还已经具备实验框架和评测基础。这意味着后续工作不需要再从头建设完整评测体系，而可以在现有实验能力之上增加控制器与基线策略的对比。就时间成本而言，这一前提极大降低了转向 controller 的风险，因为真正需要新增的重点将不再是大规模平台建设，而是问题形式化、状态定义、决策逻辑设计以及实验验证。

\subsection{为什么在当前时间节点下，Belief-State Controller 仍然可做}

虽然距离毕业设计结束时间已经不长，但 \texttt{Belief-State Controller} 方向仍然具有可行性，其原因不在于它“足够简单”，而在于它具备一个很重要的结构性优势：它可以建立在现有系统之上，以“新增控制层”的方式展开，而不要求对底层架构进行推翻式改造。

这意味着，课题的工作量可以被控制在一个较聚焦的范围内。系统的任务流程、恢复机制、人工介入逻辑和事件记录能力都已经存在，后续真正需要完成的核心工作，是围绕这些已有能力提出并验证一个统一的控制问题。只要研究范围被严格限制在“高代价科学工作流的运行时治理”，那么课题就不必再继续扩展任务覆盖面，也不必把精力投入到新的工具接入、全新流程建设或大规模产品化工作中。对于当前阶段而言，这种“在已有平台上增加新的研究层”的方式，是比继续做一个更大的助手系统更可行的。

与此同时，\texttt{Belief-State Controller} 的可行性还来自其可做“降阶版本”。前文已经讨论过，完整的 POMDP、端到端强化学习或复杂概率图推断虽然理论上更强，但在当前时间条件下风险过高。相比之下，一个结构化、显式且可解释的 \texttt{Belief-State Controller Lite} 更适合作为毕设收敛方向。它不追求在理论上求得最优控制策略，而是强调：如何定义合适的隐状态，如何根据运行时观测更新信念状态，以及如何基于更新后的信念状态做工作流级治理动作选择。这样的工作量虽然不轻，但它是可管理的；更重要的是，它仍然能保持明显的算法主线，而不会退化成纯粹的工程堆叠。

\subsection{时间可行性建立在“严格收缩研究边界”之上}

需要强调的是，controller 方向之所以可行，并不是因为它本身没有难度，而是因为它可以通过边界控制收敛到一个适合当前阶段完成的版本。若继续让课题同时承担“生命科学助手增强”“多场景扩展”“完整 benchmark 设计”“复杂控制算法求最优策略”等多重目标，那么无论选择哪个方向，都很难在不足两个月内完成。因此，本课题要成立，必须明确哪些内容是本阶段应该做的，哪些内容则应主动放弃。

首先，不应再将“构建一个更完整的生命科学助手”作为目标。其次，不应把课题范围扩展到多种不同科学场景，而应坚持以蛋白设计工作流作为唯一核心验证场景。再次，不应试图实现完整的 POMDP 求解、端到端强化学习或高度复杂的最优策略学习，而应聚焦于一个结构化、可解释、可实验验证的信念状态控制原型。最后，不应将自动 benchmark 也一并做成并列主贡献，而应将现有实验能力仅作为 controller 的验证载体。

也就是说，当前时间可行性的前提，不是“课题足够简单”，而是“课题目标被主动收束”。这种收束一旦做到，课题就会从“大而散的系统建设”转变为“围绕单一控制问题的集中验证”。

\subsection{剩余时间内最合理的工作推进逻辑}

从时间安排上看，当前阶段的工作应围绕“先收敛问题，再构造方法，最后做验证”这一顺序展开，而不能反过来先做大量系统改动，再试图从中提炼论文主线。就现实情况而言，剩余时间更适合被划分为三个连续阶段。

第一阶段是问题收敛与方法定型阶段。该阶段的核心不是编码，而是进一步压实研究对象，明确控制器所维护的状态、可用的观测、动作空间以及目标函数，并将 \texttt{Belief-State Controller} 与静态门控、规则系统或固定策略清晰区分开来。

第二阶段是控制器原型与基线对比阶段。在问题定义明确之后，需要基于现有系统形成一个可运行的控制原型，并设定至少一类静态基线作为对照。这里的关键不在于功能做得多复杂，而在于让控制器真正改变工作流治理动作，并能够在实验中体现出与固定策略的差异。

第三阶段是实验分析与论文收尾阶段。由于当前系统已有实验基础，这一阶段应重点围绕 controller 的收益和局限性展开，包括成功率、恢复代价、人工介入次数、执行时间以及失败样本分析等。

\subsection{当前方向的主要时间风险}

尽管如此，\texttt{Belief-State Controller} 方向仍存在明确的时间风险，主要集中在三个方面。第一，问题定义如果迟迟不能收敛，后续实现和实验就会持续摇摆。第二，如果在实现时试图同时追求过强的方法复杂度，例如完整的 POMDP、复杂学习式策略或多场景泛化验证，则工作量会显著失控。第三，如果实验目标设定过多，试图同时证明所有维度上的提升，那么有限时间内很难形成有说服力的核心结论。

因此，时间风险并不意味着该方向不可行，而是意味着必须坚持“单主线、单场景、单核心问题”的收敛策略。只要这一点被严格执行，本课题仍有较大概率在剩余时间内形成完整成果。

\subsection{本节结论}

综合来看，\texttt{Belief-State Controller} 方向在当前时间节点下仍然具备现实可行性，其根本原因在于：现有系统已经提供了足够完整的运行时平台和实验基础，使得后续工作可以集中在控制层的研究问题上，而无需从零搭建新的系统框架。同时，这一方向又允许通过“降阶版本”控制方法复杂度，从而避免陷入完整 POMDP 或强化学习等高风险路线。

但这种可行性并不是无条件的。它建立在三个前提之上：其一，课题必须从“生命科学助手增强”收缩为“高代价科学工作流治理控制”；其二，验证场景必须收敛到当前已有的蛋白设计工作流；其三，方法目标必须控制在一个可解释、可验证的 \texttt{Belief-State Controller Lite} 范围内。
